\subsection{Функции}

Функция — это именованный блок кода, который можно вызывать многократно. Функции делают код чище и понятнее.

Ключевое слово \texttt{def} объявляет функцию. \texttt{return} возвращает результат.

\begin{codefile}[python]{python/01_python_basics/07_functions.py}
\end{codefile}

\begin{infobox}[Область видимости]
Переменные, созданные внутри функции, не видны снаружи. Это называется локальной областью видимости.
\end{infobox}

\textbf{Параметр} --- имя в скобках после имени функции: при вызове \texttt{calculate\_area(5)} в нём оказывается значение 5. Одна функция с параметром заменяет несколько почти одинаковых кусков кода. Функция без \texttt{return} тоже полезна --- она просто выполняет действия, например печатает или включает подсветку; результат её вызова --- \texttt{None}.

\subsubsection*{Задачи для самостоятельного решения}

\begin{enumerate}
    \item Напишите функцию \texttt{say\_hello()}, которая печатает "Привет!". Вызовите её.
    \item Напишите функцию \texttt{greet(name)}, которая печатает "Привет, <name>!".
    \item Напишите функцию \texttt{add(a, b)}, которая возвращает сумму двух чисел.
    \item Напишите функцию \texttt{square(x)}, которая возвращает x в квадрате.
    \item Напишите функцию \texttt{my\_max(a, b)}, которая возвращает большее из двух чисел.
    \item Напишите функцию \texttt{is\_even(n)}, которая возвращает \texttt{True}, если число четное.
    \item Напишите функцию \texttt{perimeter(a, b)}, возвращающую периметр прямоугольника.
    \item Напишите функцию \texttt{m\_to\_cm(meters)}, которая переводит метры в сантиметры.
    \item Напишите функцию \texttt{check\_battery(voltage)}. Если < 3.5 -> "Low", иначе "OK".
    \item Напишите функцию \texttt{my\_abs(x)}. Если x < 0, возвращает -x, иначе x.
    \item Напишите функцию \texttt{avg(a, b)}, возвращающую среднее арифметическое.
    \item Напишите функцию \texttt{log\_error(msg)}, которая печатает "ERROR: <msg>".
    \item Напишите функцию \texttt{repeat(s, n)}, возвращающую строку s, повторенную n раз.
    \item Напишите функцию \texttt{flight\_time(cap, cons)}, считающую время полета (емкость / расход).
    \item Напишите функцию \texttt{full\_name(first, last)}, возвращающую "Фамилия Имя".
\end{enumerate}

\begin{taskbox}[Конвертер]
Напишите функцию \texttt{celsius\_to\_fahrenheit(c)}, которая возвращает температуру по Фаренгейту по формуле $F = C \times 1.8 + 32$. Вызовите её в цикле для 0, 100 и $-40$ и выведите пары «Цельсий --- Фаренгейт».
\expected{\texttt{0 32.0}, \texttt{100 212.0}, \texttt{-40 -40.0}. Решение --- в конце файла \texttt{examples/python-answers/07\_functions\_answers.py}.}
\mission{Python}{25--26}
\end{taskbox}
